\section{Fourier bounds for one marginal, uniformly in the index}\label{j:fourier:section}
Two transformations are needed. A compensated inversion formula bounds the
smooth-payoff error $E_n$ in conservation; Prawitz's inequality bounds the
CDF itself. Both use characteristic functions of the same marginal. The
formulas below are pointwise inequalities with explicit moment and support
hypotheses, so they also cover triangular arrays of changing laws. The finite
certificates evaluate these formulas on closed cells; they do not supply an
asymptotic or infinite-domain analytic argument.

\begin{lemma}[Gaussian normalizations]\label{j:fourier:constants}
One has $\|\phi''\|_\infty=\phi(0)<2/5$,
$\|\phi'''\|_\infty\le2/\pi$, and
\begin{gather*}
 \frac{333}{106}<\pi<\frac{22}{7},\quad
 \sqrt\pi>\frac{4431}{2500},\quad
 \sqrt{\pi/2}<\frac{251}{200},\\
 \phi(0)>0.39886>\frac{997}{2500},\quad
 \frac{819}{2000}<\CE<\frac{20487}{50000}.
\end{gather*}
\end{lemma}
\begin{proof}
Differentiation and integration by parts in the Gaussian integral give
$\phi(x)=(2\pi)^{-1}\int_\R e^{-t^2/2-itx}dt$; its normalization follows
by squaring the integral and using polar coordinates. Absolute convergence gives
\[
 \|\phi^{(k)}\|_\infty\le\frac1{2\pi}\int_\R |t|^ke^{-t^2/2}dt.
\]
This proves the third derivative bound. For the second derivative,
$(x^2-1)e^{-x^2/2}$ has absolute maximum one (the other candidate is
$2e^{-3/2}<1$).
For the rational constants put
$s_k(x)=\sum_{j=0}^{k-1}(-1)^jx^{2j+1}/(2j+1)$.
Tangent addition, with $0<4\arctan(1/5)-\arctan(1/239)<\pi/2$,
gives Machin's identity $\pi=16\arctan(1/5)-4\arctan(1/239)$.
The geometric remainder in $1/(1+t^2)$, integrated from zero, gives
\[
 \pi_-=16s_6(1/5)-4s_1(1/239)<\pi
 <16s_5(1/5)-4s_2(1/239)=\pi_+.
\]
Exact rational comparison gives $333/106<\pi_-<\pi_+<22/7$ and
$\pi_->(4431/2500)^2$, $\pi_+<2(251/200)^2$.
Also $(0.39886)^2<7/44$.
Finally, for $r_-=3162277660/10^9$, $r_+=3162277661/10^9$,
$r_-^2<10<r_+^2$ and
\[
 (819/2000)^2<\frac{19+6r_-}{72\pi_+}
 <\CE^2<\frac{19+6r_+}{72\pi_-}<(20487/50000)^2.
\]
All denominators are positive; the displayed finite sums specify the
comparisons as integer cross-multiplications.
\end{proof}

For $n\ge2$, write
\[
 f(t)=\E e^{itX},\quad e_n(t)=e^{-nt^2/2}(1-in\gamma t^3/6),\quad
 V_n(t)=f(t)^n-e_n(t),
\]
\[
 I_0(H)=\int_0^{1/2}\frac{|\Im(e^{-itH}V_n(t))|}{t^3}\,dt,\qquad
 I_\infty=\int_{1/2}^\infty\frac{|V_n(t)|}{t^3}\,dt.
\]
\subsection{Compensated inversion, including atoms}

The compensation below makes absolute Fubini valid, including at atoms.

\begin{lemma}\label{j:fourier:inversion}
If a finite signed measure $\mu$ has mass and first moment zero and
$\int x^2\,d|\mu|(x)<\infty$, then, for every $H\in\mathbb R$,
\[
 \int(x-H)|x-H|\,d\mu(x)
 =-\frac4\pi\int_0^\infty
 \frac{\Im(e^{-itH}\widehat\mu(t))}{t^3}\,dt.
\]
The integral on the right is absolutely convergent.
\end{lemma}
\begin{proof}
For real $y$, substitution and integration by parts give
\[
 \int_0^\infty\frac{ty-\sin(ty)}{t^3}\,dt=\frac\pi4 y|y|.
\]
The boundary terms vanish: near zero the numerator is $O(t^3)$, and
at infinity it is $O(t)$. The integrand has the sign of $y$, since
$v-\sin v\ge0$ for $v\ge0$. To calculate the constant, nonnegative
Tonelli and elementary Laplace integration give
\[
 \int_0^\infty\frac{1-\cos t}{t^2}\,dt
 =\int_0^\infty s\left(\frac1s-\frac{s}{s^2+1}\right)ds
 =\frac\pi2.
\]
The absolute double integral of the compensated expression against
$d|\mu|(x)$ equals
\[
 \frac\pi4\int(x-H)^2d|\mu|(x)<\infty.
\]
Absolute Fubini is therefore valid. The linear compensation disappears
because $\int(x-H)d\mu=0$, proving both assertions. No continuity of a
CDF at $H$ is involved.
\end{proof}

Apply this to the sum law minus the scale-$\sqrt n$ pushforward of
the signed density
\[
 \phi(x)\{1+\gamma(x^3-3x)/(6\sqrt n)\}.
\]
This signed density has mass one, first moment zero and finite absolute
second moment. Its transform after scaling is $e_n$. Gaussian integration
by parts, using $q_z(x)=(x-z)|x-z|$ and $q_z'''=4\delta_z$, yields
\[
 \varepsilon_n:=\E q_{H/\sqrt n}(S_n/\sqrt n)
 -\E q_{H/\sqrt n}(Z)-\frac{2\gamma\phi(H/\sqrt n)}{3\sqrt n}
 =-\frac4{\pi n}\int_0^\infty
 \frac{\Im(e^{-itH}V_n(t))}{t^3}\,dt.
\]
Consequently the error $E_n=-n^{3/2}\varepsilon_n/2$ satisfies
\begin{equation}\label{j:fourier:ledger}
 |E_n|\le\frac{2\sqrt n}{\pi}\{I_0(H)+I_\infty\}.
\end{equation}
The identity and bound hold for each $n$ and $H$ separately.

\subsection{Characteristic-function bounds from common moments}
Let $X$ be centered with $\E X^2=1$, $\rho=\E|X|^3<\infty$,
$\gamma=\E X^3$, and $f(u)=\E e^{iuX}$. In this section put
\[
 \begin{gathered}
 h=\E|X|,\quad m=\E X_+=\E X_-=h/2,\quad
 A=\E X_+^2,\quad B=\E X_-^2,\\
 \zeta=2AB/m,\qquad b(r)=\frac{\sqrt{r^2+8}-r}{2}.
 \end{gathered}
\]
Here $X_-=\max\{-X,0\}$, so $A+B=1$ and all of $h,m,A,B$ are
positive. The letter $b$ is used for this root function to avoid confusion
with the signed second moment $\beta=\E[X|X|]$.

\begin{lemma}[Anchored characteristic function]\label{j:fourier:anchor}
For $0\le u\le\pi$,
\[
 \cos u\le\Re f(u)\le\cos u+(\rho-1)u^3/6,
 \qquad
 |\Im f(u)|^2\le\frac{u^6}{36}(\rho-1)(\rho+5/3).
\]
\end{lemma}
\begin{proof}
Write $\operatorname{sinc}t=\int_0^1\cos(st)\,ds$. For real $a,b$,
\[
 |\operatorname{sinc}a-\operatorname{sinc}b|
 \le |a-b|/2,\qquad
 |\operatorname{sinc}a-\operatorname{sinc}b|
 \le |a^2-b^2|/6.
\]
The first follows from the Lipschitz bound for cosine. For the second,
$v\mapsto\cos\sqrt v$ has derivative of absolute value at most $1/2$
on $v\ge0$, including its limit at zero; integrate with respect to $s$.
For $r\ge0$, integration from $1$ to $r$ of the first estimate gives
\[
 \left|\cos(ur)-\cos u+\frac{u\sin u}{2}(r^2-1)\right|
 \le \frac{u^3}{12}(r-1)^2(2r+1).
\]
Since $\E(|X|-1)^2(2|X|+1)=2(\rho-1)$, averaging proves the
upper real bound. For the lower bound, when $0<u\le\pi$ set
$a=\sin u/(2u)$ and $J(v)=\cos v+av^2$. The function
$\operatorname{sinc}v$ decreases on $[0,\pi]$, so $J$ decreases up to
$u$ and increases from $u$ to $\pi$. For $v\ge\pi$,
$J(v)\ge-1+a\pi^2=J(\pi)\ge J(u)$. Evenness gives the same conclusion
for negative $v$. Average $J(uX)\ge J(u)$ and use $\E X^2=1$.

The second sinc estimate, together with $\E X=0$, gives
$|\Im f(u)|\le u^3\E[|X|\,||X|^2-1|]/6$.
Put $W(r)=(r-1)^2(2r+1)$. Cauchy's inequality applies to the factorization
$r|r^2-1|=\sqrt{W(r)}\,r(r+1)/\sqrt{2r+1}$.
The exact identity
\[
 r^3+\frac{29r^2+1}{18}
 -\frac{2r^2(r+1)^2}{2r+1}
 =\frac{(r-1)^2(4r+1)}{18(2r+1)}\ge0
\]
bounds the second squared factor in expectation by $(\rho+5/3)/2$;
the first has expectation $2(\rho-1)$. All expectations are finite under
the stated third-moment hypothesis. The value at $u=0$ is immediate.
\end{proof}

\begin{lemma}[One common moment surrogate]\label{j:fourier:surrogate}
For an independent copy $X'$, put
$\eta=\E|X-X'|^3/2-\rho$. Then
\[
 \eta\le h\le1,\qquad \zeta\ge h,\qquad
 \rho+\zeta\ge2/h,\qquad \zeta\ge b(\rho),
 \qquad \gamma^2+\zeta^2\le\rho^2.
\]
For every real $u$ one also has
$36|\Im f(u)|^2\le u^6(\rho^2-\zeta^2)$.
If $|X|\le R$ almost surely, then additionally
\[
 \eta\le b(\rho)+\frac{(R-\sqrt{\rho^2+8})_+^2}{4R}.
\]
In particular $R\le\sqrt{\rho^2+8}$ implies $\eta\le b(\rho)$.
\end{lemma}
\begin{proof}
The scalar inequality
$|x-y|^3\le(x-y)^2(|x|+|y|)$, independence and centering give
$\E|X-X'|^3\le2\rho+2h$; Cauchy gives $h\le1$.
Signwise Cauchy gives
\[
 A^2\le m\E X_+^3,\quad B^2\le m\E X_-^3,
 \quad \frac{m^2}{A}+\frac{m^2}{B}\le\Pp(X\ne0)\le1.
\]
Thus $AB\ge m^2$, $\zeta\ge2m=h$, and
$\rho+\zeta\ge(A+B)^2/m=2/h$. Multiplication yields
$\zeta(\rho+\zeta)\ge2$, proving $\zeta\ge b(\rho)$.

For completeness, the sine-circle bound uses an exact mixture, not an
extreme-point assertion about convolution. Let $\mu_+$ and $\mu_-$ be the
restrictions of the law of $|X|$ to its positive and negative sign classes.
On $a,b>0$ define
\[
 d\nu(a,b)=\frac{a+b}{m}\,d\mu_+(a)d\mu_-(b).
\]
Integrating the centered law with atoms $a,-b$ and probabilities
$b/(a+b),a/(a+b)$ against $\nu$ recovers the nonzero part of the original
law; the zero atom is added separately. For this two-point law the second
sinc estimate implies
\[
 \frac{6|\Im f_{a,b}(u)|}{|u|^3}
 \le\frac{ab|a^2-b^2|}{a+b}.
\]
Together with $e_{a,b}=2a^2b^2/(a+b)$, the right side is one coordinate
of a vector of length $ab(a^2+b^2)/(a+b)$. Minkowski's inequality in
$\R^2$ therefore bounds the vector
$(\zeta,6\Im f(u)/u^3)$ by $\rho$. The third-moment assumption
justifies all integrals. Letting $u\to0$ proves the assertion for $\gamma$:
Taylor's remainder is dominated by $|X|^3/6$, so dominated convergence
applies after using $\E X=0$.

For the final assertion the exact same-sign deficit is
\[
 h-\eta=\E\!\left[
 \min(|X|,|X'|)(X-X')^2\ind_{\{XX'>0\}}\right].
\]
On a common sign class, $\min(a,b)\ge ab/R$. Expansion of the two
independent class integrals gives
\[
 h-\eta\ge\frac{h\rho-2(A^2+B^2)}{R}
 =\frac{h(\rho+\zeta)-2}{R}.
\]
If $h\le b(\rho)$, the already proved $\eta\le h$ suffices. Otherwise,
write $b=b(\rho)$, use $\zeta\ge h$ and $b(\rho+b)=2$ to obtain
\[
 \eta-b\le\frac{(h-b)(R-\rho-h-b)}{R}\le0,
\]
when $R\le\rho+2b=\sqrt{\rho^2+8}$. For general $R$, the quadratic
$y(R-\sqrt{\rho^2+8}-y)$ on $y\ge0$ is at most
$(R-\sqrt{\rho^2+8})_+^2/4$. This proves the stated remainder too.
In particular radius three suffices for every $\rho\ge1$.
\end{proof}

For a universal calculation put $H_* =\max\{h,b(\rho)\}$.
Then $b(\rho)\le H_*\le1$, $\eta\le H_*$, $\zeta\ge H_*$ and
$\rho+\zeta\ge2/H_*$. On a closed cell
$\rho\in[\ell,r]$, $H_*\in[a,b]$, it is consequently safe to use
\[
 \eta\le b,\qquad
 \zeta\ge\max\{a,2/b-r,b(r)\}.
\]
Every cell uses both bounds from the same marginal law. A lower bound on
$H_*$ is not a lower bound on $h$; no later sign-probability argument may
substitute it for one. If support $|X|\le R\le\sqrt{\ell^2+8}$ has independently been acquired,
the last part of Lemma~\ref{j:fourier:surrogate} instead permits
$\eta\le b(\ell)$ and $\zeta\ge b(r)$. For larger $R$ its displayed
nonnegative remainder must be retained.

\subsection{Keeping the cubic compensation in the high frequencies}
\label{j:fourier:compensated}
Suppose now that $|X|\le R$ and $\rho\le r$. We may use
\begin{align*}
 M_4&=\min\left\{Rr,1+\frac{(R+1)^2}{R+1/2}(r-1)\right\},\\
 M_5&=\min\left\{RM_4,
 1+\frac{R^3+2R^2+3R+3/2}{R+1/2}(r-1)\right\}
\end{align*}
as upper bounds on $\E X^4$ and $\E|X|^5$ of the same law. To verify the
less immediate branches, set $y=|X|$ and expand
\[
 \rho-1=\E(y-1)^2(y+1/2),\quad
 \E X^4-1=\E(y-1)^2(y+1)^2,
\]
\[
 \E|X|^5-1=\E(y-1)^2(y^3+2y^2+3y+3/2).
\]
The ratios of the last two factors to $y+1/2$ increase on $y\ge0$:
their derivatives are respectively
$y(y+1)/(y+1/2)^2$ and $y(4y^2+7y+4)/(2(y+1/2)^2)$.
The other branches follow from $y\le R$. These moment bounds require
no bound on the product-gap defect.

Put $G=\min\{(r-1)(r+5/3),r^2-Z^2\}$, where $Z>0$ is a proved
lower bound for $\zeta$ on the whole cell; $Z=b(r)$ always suffices.
Then $\gamma^2\le G$: the first branch follows by letting $u\to0$ in
Lemma~\ref{j:fourier:anchor}, and the second from
Lemma~\ref{j:fourier:surrogate}.
Assume that a proved nonnegative envelope $m(u)$ bounds $|f(u)|$ on
the interval being integrated. Set
\[
 g=e^{-u^2/2},\quad v=-i\gamma u^3/6,\quad
 b_u=\sqrt G\,u^3/6,\quad W=g\sqrt{1+b_u^2},
\]
and, for $0\le u\le\pi$,
\[
 \begin{gathered}
 A_u=\max\{|\cos u-g|,
 |\cos u+(r-1)u^3/6-g|\},\\
 B_u=\min\{M_4u^4/24, M_5u^5/120\}+b_u(1-g).
 \end{gathered}
\]
Taylor's theorem gives
\[
 |\Im f(u)+\gamma u^3/6|\le\min\{M_4u^4/24,M_5u^5/120\}.
\]
The anchored real estimate therefore gives
$|f-g(1+v)|\le\sqrt{A_u^2+B_u^2}$. Telescoping the $n$th powers and
then applying the integral Taylor remainder to $(1+sv)^n$, $0\le s\le1$,
proves, for every integer $n\ge2$,
\begin{align*}
 |f(u)^n-g^n(1+nv)|
 &\le\sqrt{A_u^2+B_u^2}\sum_{j=0}^{n-1}m(u)^jW^{n-1-j}
 +\binom n2 b_u^2g^n(1+b_u^2)^{(n-2)/2},\\
 |f(u)^n-g^n(1+nv)|
 &\le m(u)^n+g^n\sqrt{1+n^2b_u^2}.
\end{align*}
One may take their pointwise minimum. These formulas use no quotient of
$m(u)-W$, so coincident and zero values cause no exceptional cases.

On a closed frequency cell $[a,b]\subset[1/2,T]$, a rigorous upper
bound for this minimum is multiplied by the exact weight
$\int_a^b u^{-3}\,du=(a^{-2}-b^{-2})/2$. The infinite remainder is
bounded without a periodicity assumption: for $T>0$,
\[
 \int_T^\infty\frac{|f(u)^n-g(u)^n(1+nv(u))|}{u^3}\,du
 \le\frac1{2T^2}
 +\frac{e^{-nT^2/2}}{nT}
 \left(T^{-3}+\frac{n\sqrt G}{6}\right).
\]
Indeed $|f|\le1$, $\sqrt{1+x^2}\le1+|x|$, $u^{-3}\le T^{-3}$
on the normal part, and
$\int_T^\infty e^{-nu^2/2}\,du\le e^{-nT^2/2}/(nT)$.

To preserve cancellation in the real residual on a closed frequency cell,
put $d_0(u)=\cos u-e^{-u^2/2}$ and
$d_1(u)=d_0(u)+(r-1)u^3/6$. Their derivatives are
\[
 d_0'(u)=-\sin u+ue^{-u^2/2},\qquad
 d_1'(u)=d_0'(u)+(r-1)u^2/2.
\]
If $D_j$ encloses $d_j'$ throughout a cell of width $w$ and midpoint
$u_0$, the fundamental theorem of calculus gives
$d_j(u)\in d_j(u_0)+[-w/2,w/2]D_j$.
Intersecting this with a direct interval evaluation remains valid.

\subsection{The compatible modulus profile}
\label{j:fourier:profile}
The scalar certificate uses $\kappa=62/625$ and proves
\begin{equation}\label{j:fourier:scalar-cosine}
 \cos x\le1-x^2/2+\kappa|x|^3\qquad(x\in\R).
\end{equation}
Here is the entire reduction to its finite check. For $x\ge625/124$ the
right side is at least one. On the remaining interval, the degree-sixteen
Taylor polynomial $P_{16}$ bounds cosine above: the sixteenth derivative
of $\cos-P_{16}$ is $\cos-1\le0$, and its lower derivatives vanish at
zero. The degree-thirteen polynomial
\[
 Q(x)=\kappa+\sum_{j=2}^{8}\frac{(-1)^{j+1}x^{2j-3}}{(2j)!}
\]
satisfies $x^3Q(x)=1-x^2/2+\kappa x^3-P_{16}(x)$.
On the eighty cells $[i/16,(i+1)/16]$, $0\le i<80$, and the last
cell $[5,625/124]$, the exact verifier checks the inverse Bernstein
identity and all fourteen coefficients. The smallest coefficient is
$2479189/65585520000>0$. Lemma~\ref{j:cert:bernstein} proves positivity
on every closed cell; evenness completes \eqref{j:fourier:scalar-cosine}.

Put $a=1/2-4\kappa=129/1250$, $C=1-\cos4$ and $A_0=16a/C$.
The profile
\[
 q(z)=\begin{cases}
 1/2-\kappa z,&0\le z\le4,\\
 A_0(1-\cos z)/z^2,&4\le z\le2\pi,\\
 0,&z\ge2\pi
 \end{cases}
\]
is nonnegative, convex and nonincreasing on $[0,\infty)$, and satisfies
$\cos z\le1-z^2q(z)$. Indeed \eqref{j:fourier:scalar-cosine} gives
$0<A_0\le1$. The pieces agree at their endpoints. On the middle interval
the first derivative is nonpositive and the second has the sign of
\[
 B_0(z)=6+(z^2-6)\cos z-4z\sin z.
\]
On $[4,3\pi/2]$, $B_0'(z)=-2z\cos z+(2-z^2)\sin z\ge0$ and
$B_0(4)\ge4$: indeed $7\pi/6\le4\le3\pi/2$ implies
$\sin4\le-1/2$, and $\cos4\ge-1$, so
$6+10\cos4-16\sin4\ge4$. On $[3\pi/2,2\pi]$, $B_0(z)\ge6$.
The derivative at $2\pi$ is zero. The derivative jump at $4$ has the
correct sign because $258\sin4+119(1-\cos4)>0$. Its exact lower bound is
obtained by replacing sine by its degree-fifteen Taylor lower bound and
cosine by $P_{16}$. The same scalar checker verifies this rational
inequality and the needed $25/8<\pi<22/7$ by Machin's formula.

To apply the profile, weight the distribution of $X-X'$ by
$(X-X')^2/2$. This is a probability measure whose first absolute moment
is $\rho+\eta$. Jensen's inequality gives, for $u\ge0$,
\[
 |f(u)|^2=\E\cos(u(X-X'))
 \le1-2u^2q(u(\rho+\eta)).
\]
If $\rho+\eta\le \mathcal D$, monotonicity permits $\mathcal D$ on the right. Thus the
linear part gives $1-u^2+2\kappa \mathcal Du^3$; when $4\le \mathcal Du\le2\pi$ the
curved part gives $1-2A_0(1-\cos(\mathcal Du))/\mathcal D^2$.
Take the minimum with one and with the coordinate-disk estimate from
Lemmas~\ref{j:fourier:anchor}--\ref{j:fourier:surrogate}. The upper
square of a real interval is the larger of the squares of its two
endpoints, even when that interval crosses zero or lies below it.

\subsection{A Fourier majorant of a closed half-line}

\begin{lemma}\label{j:fourier:closed-majorant}
For every integrable real random variable $Y$, every $h\in\R$ and $A>0$,
\begin{equation}\label{j:fourier:closed-majorant-formula}
 \Pp(Y\le h)\le\frac12+
 2\lim_{\epsilon\downarrow0}\Re\int_\epsilon^1
 K(t)e^{-iAth}\E e^{iAtY}\,dt,
\end{equation}
where
\[
 K(t)=\frac{1-t}{2}+\frac i2
 \left((1-t)\cot(\pi t)+\frac1\pi\right).
\]
The limit concerns the real part; the complex integral itself may diverge.
\end{lemma}
\begin{proof}
Put $W(x)=(\sin(\pi x)/(\pi x))^2$, with $W(0)=1$, and
\[
 H_B(x)=\sum_{m\ge1}\{W(x-m)-W(x+m)\}+2xW(x).
\]
The series converges absolutely and locally uniformly. It is odd, vanishes
at zero, and takes value $\operatorname{sgn}(j)$ at every nonzero integer.
The partial-fraction identity
\[
 \sum_{m\in\mathbb Z}(x-m)^{-2}=\pi^2\csc^2(\pi x)
 \quad(x\notin\mathbb Z)
\]
follows by integrating $\pi\cot(\pi z)/(z-x)^2$ around the square
with vertices $\pm(N+1/2)\pm i(N+1/2)$: the boundary integral is
$O(N^{-1})$, and the residues are $(m-x)^{-2}$ and
$-\pi^2\csc^2(\pi x)$. For positive noninteger $x$ this identity gives
\[
 H_B(x)-1=\left(\frac{\sin\pi x}{\pi}\right)^2
 \left(\frac2x-\frac1{x^2}-2\sum_{m\ge1}(x+m)^{-2}\right).
\]
Integral comparison of the decreasing function $(x+s)^{-2}$ yields
\[
 \sum_{m\ge1}(x+m)^{-2}\le\frac1x
 \le\frac1{x^2}+\sum_{m\ge1}(x+m)^{-2}.
\]
Thus $|H_B(x)-\operatorname{sgn}(x)|\le W(x)$; oddness and the integer
values supply the remaining points. Consequently
\[
 M(x):=\frac{1-H_B(x)+W(x)}2\ge\ind_{\{x\le0\}},
 \qquad M(0)=1,\qquad 0\le M(x)-\ind_{\{x\le0\}}\le W(x).
\]
Elementary integration gives
\[
 W(x)=2\int_0^1(1-t)\cos(2\pi tx)\,dt,\qquad
 2xW(x)=\frac2\pi\int_0^1\sin(2\pi tx)\,dt.
\]
Inserting the first formula into the shifted series with Abel weights
$r^m$, $0<r<1$, gives
\[
 \sum_{m\ge1}r^m\{W(x-m)-W(x+m)\}
 =4\int_0^1(1-t)\sin(2\pi tx)
 \frac{r\sin(2\pi t)}{1-2r\cos(2\pi t)+r^2}\,dt.
\]
The fraction tends to $\tfrac12\cot(\pi t)$ and is bounded in absolute
value by $\tfrac12|\cot(\pi t)|$, since its denominator is
$(1-r)^2+4r\sin^2(\pi t)$. The resulting majorant is integrable:
the sine vanishes at zero and $(1-t)\cot(\pi t)$ is bounded at one.
Dominated convergence therefore proves
\[
 H_B(x)=2\int_0^1
 \left((1-t)\cot(\pi t)+\frac1\pi\right)\sin(2\pi tx)\,dt.
\]
Combining the formulas yields
\[
 M(x)=\frac12+2\lim_{\epsilon\downarrow0}\Re\int_\epsilon^1
 K(t)e^{2\pi itx}\,dt
 =1+2\Re\int_0^1K(t)(e^{2\pi itx}-1)\,dt.
\]
The last integral is absolutely convergent. For integrable $U$,
$\E|e^{2\pi itU}-1|\le2\pi t\E|U|$; together with $K(t)=O(t^{-1})$
at zero and boundedness at one, this permits absolute Fubini.
Apply the identity to $U=A(Y-h)/(2\pi)$ and take expectations.
The value $M(0)=1$ proves the claimed closed-event convention even at atoms.
\end{proof}

For $Y=(X_1+\cdots+X_4)/2$ and $A=6$, independence gives
$\E e^{6itY}=f(3t)^4$. Thus the lemma supplies exactly the
factor six, positive imaginary kernel, and boundary value used in the
four-summand comparison below.

\subsection{From characteristic functions to the CDF}
\label{j:fourier:prawitz}
We use Prawitz's smoothing inequality in the form recorded in
\cite[Lemma 4.1, p.~263]{Shevtsova2012}. For $n\ge2$, $c>0$, $0<\tau\le1$, put
\[
 K_2(t)=(1-t)+i\{(1-t)\cot(\pi t)+1/\pi\},\qquad 0<t<1.
\]
With $T=c\sqrt n$, that inequality gives
\begin{align*}
 \Delta_n\le{}&\int_0^\tau |K_2(t)|
   |f(ct)^n-e^{-nc^2t^2/2}|\,dt
 +\int_\tau^1|K_2(t)|\,|f(ct)|^n\,dt\\
 &+\int_0^\tau|K_2(t)-i/(\pi t)|e^{-nc^2t^2/2}\,dt
 +\frac1\pi\int_\tau^\infty e^{-nc^2t^2/2}\frac{dt}{t}.
\end{align*}
Our standardized iid laws satisfy its independence, centering and positive
total-variance hypotheses. The cited convention uses a strict CDF.
Continuity of $\Phi$ makes the supremum identical for closed CDFs: take
one-sided limits at each evaluation point. Thus atoms require no smoothing
of the original law.

For $g(u)=e^{-u^2/2}$ and envelopes $m\ge|f|$, $E_0\ge|f-g|$,
the first integrand uses
$|f^n-g^n|\le\min\{E_0\sum_{j=0}^{n-1}m^jg^{n-1-j},m^n+g^n\}$.
Every literal integer $n$ is evaluated separately. The original $20$--$24$ integral verifier
covers $[\epsilon,\tau]$ and $[\tau,1-\epsilon]$ by closed cells,
where $\epsilon=1/1024$; its numerical parameter domain is
$0<c<\pi$ and $\epsilon<\tau\le3/10<1-\epsilon$.
In particular $c\tau<1$ licenses the stable low-frequency Taylor bounds.
On a cell of width $w$ and midpoint $t_0$,
it uses either $w\sup F$ or
$w\{F(t_0)+w^2\max(0,\sup F'')/24\}$ for a fixed smooth majorant
$F$. The latter follows by Taylor expansion and integration; a failed
derivative-domain guard triggers the former bound on the entire cell.
Switching branches at a sampled midpoint would not justify this rule.

For the omitted endpoint pieces, the respective upper bounds are
\begin{align*}
 L_\epsilon&=n\left[
 \frac{rc^3}{6}\left(\frac{\epsilon^4}{4}
       +\frac{\epsilon^3}{3\pi}\right)
 +\frac{c^4}{8}\left(\frac{\epsilon^5}{5}
       +\frac{\epsilon^4}{4\pi}\right)\right],\\
 N_\epsilon&=\epsilon+\pi\epsilon^2/2,\qquad
 H_\epsilon=\epsilon^2/2+\pi\epsilon^3/3.
\end{align*}
They follow from $|K_2(t)|\le1+1/(\pi t)$ near zero, cubic Taylor
error for $f$, $e^{-u^2/2}-1+u^2/2\le u^4/8$, and
\[
 0\le1/x-\cot x\le\frac{x}{3(1-x^2/6)}\le x\quad(0<x\le1).
\]
For this last estimate use
$\sin x-x\cos x=\int_0^x t\sin t\,dt\le x^3/3$ and
$\sin x\ge x-x^3/6$. Apply the same identity with $1-t$ near one.
Finally the infinite Gaussian term is at most
$e^{-nc^2\tau^2/2}/(\pi nc^2\tau^2)$, by
$\int_a^\infty e^{-v}v^{-1}\,dv\le e^{-a}/a$.
For a strip $\rho\in[\ell,r]$, multiply the full four-term bound by
$\sqrt n/\ell$, with outward enclosures. A complete closed strip cover
and a strict rational comparison then establish the claimed CDF bound.


\subsection{Sharper scalar bounds for the same integrals}
\label{j:fourier:quadrature}
For $a>0$ write $E_1(a)=\int_a^\infty e^{-v}v^{-1}\,dv$.
Integration by parts, followed by Jensen's inequality for the convex
function $t\mapsto(a+t)^{-2}$ and an exponential random variable of mean
one, gives
\[
 E_1(a)=\frac{e^{-a}}a-e^{-a}\int_0^\infty
                \frac{e^{-t}}{(a+t)^2}\,dt
 \le e^{-a}\left(\frac1a-\frac1{(a+1)^2}\right).
\]
The last Gaussian term in Prawitz's inequality is $E_1(a)/(2\pi)$,
where $a=nc^2\tau^2/2$. This improvement has no additional moment or
support hypothesis.

Some small-index certificates require $\tau>3/10$. Their low-frequency
range is still controlled explicitly. Put $U=6/5$. On $0\le x\le U$,
the alternating sine bounds and integration give
\[
 \sin x\ge x\left(1-x^2/6+x^4/120-x^6/5040\right),\qquad
 0\le\sin x-x\cos x\le x^3/3-x^5/30+x^7/840.
\]
The denominator below is at least $1-U^2/6-U^6/5040>0$, and the
numerator factor is at least $1/3-U^2/30>0$. Consequently
\[
 0\le1/x-\cot x\le L(x):=
 \frac{x(1/3-x^2/30+x^4/840)}
      {1-x^2/6+x^4/120-x^6/5040},
 \qquad L(0)=0.
\]
It follows that
\[
 |K_2(t)-i/(\pi t)|
 \le(1-t)\sqrt{1+L(\pi t)^2}
 \quad(0<t\le19/50).
\]
At the high endpoint put $s=1-t$. The identity
$K_2(t)=s+i\,s(1/(\pi s)-\cot(\pi s))$ gives
$|K_2(t)|\le s\sqrt{1+L(\pi s)^2}$ for $t\ge3/4$.
The rational upper bound $\pi<22/7$ implies
$\pi(19/50)<6/5$ and $\pi/4<6/5$.
The numerical quadrature must therefore declare
$\epsilon<\tau\le19/50<1-\epsilon$, $c<\pi$, and $c\tau\le6/5$.
It differentiates the displayed smooth rational majorant on whole cells;
it does not differentiate an interval-selected minimum.

The cancellation series on this larger range also have explicit tails.
Truncate the even series for cosine and $e^{-u^2/2}$ at degree twenty.
For their difference, the error in each $t$ derivative of order
$j=0,1,2$, with $u=ct\in[0,U]$, is at most
\[
 \frac{2(22)_j c^j u_{\max}^{22-j}}{2^{11}11!}.
\]
Indeed, for terms of Gaussian index $k\ge11$, the ratios of consecutive
absolute derivative terms are bounded respectively by
\[
 \frac{U^2}{2(k+1)},\qquad \frac{U^2}{2k},\qquad
 \frac{U^2(2k+1)}{2k(2k-1)}.
\]
All decrease with $k$ and are less than $1/8$ at $k=11$.
The cosine ratios are
$U^2/\{(2k+2-j)(2k+1-j)\}<1/8$.
The first cosine tail is $2^{11}11!/22!<1/4$ times the first Gaussian
tail, for every $j$. Summing both geometric tails gives a factor less
than $(1+1/4)/(1-1/8)=10/7<2$.
Uniform absolute convergence justifies these differentiations.
For the removable function $2(\cos u-e^{-u^2/2})/u^3$, divide the series
symbolically, so its value tail is bounded by
$4u_{\max}^{19}/(2^{11}11!)$ and its value at zero is obtained directly
from the series. No numerical division by a nearly zero frequency is
needed.

\subsection{Uniform low-frequency remainder}

The next estimate requires actual fourth through sixth moments, acquired
from support or separately bounded. A third-moment ceiling alone is insufficient.

Assume $n\ge N>2$ and common bounds
\[
 \mu_4:=\E X^4\le M\le111/10,\quad |\gamma|\le g\le37/20,
 \quad\mu_6:=\E X^6\le M_6,
\]
and put
\[
 \mathfrak a=(3-\mu_4)/24,\quad \mathfrak z=(10\gamma-\E X^5)/120,
 \quad |\mathfrak a|\le A,\quad |\mathfrak z|\le Z.
\]
Here $A,Z,M_6$ are nonnegative common ceilings, not quantities inferred
from bounded third moments. Define
\[
 E=M_6/720+g^2/72+(1/2+g^2/144)^3/5,\qquad
 B=A+Z/2+E/4,
\]
\[
 a_0=\frac{N-1}{2N},\quad b_0=\frac{45-M}{96}(1-2/N)>0,
 \quad\lambda_N=\sqrt{\pi/2}\left(\frac N{N-1}\right)^{3/2}.
\]
\begin{proposition}\label{j:fourier:low}
For all the laws and thresholds above,
\[
 I_0(H)\le\frac{L_0+L_H|H|}{\sqrt n}+\frac{L_2}{n},
\]
where, for any $c_0\ge\lambda_N$ and $c_1\ge\lambda_N/2$,
\begin{align*}
 L_0&=c_0Z+c_1Ag+251g^3/17280,\\
 L_H&=c_0A+251g^2/4800,\\
 L_2&=E/(2a_0^2)+B^2/(2b_0^3).
\end{align*}
\end{proposition}
\begin{proof}
Let $h(t)=\exp(-t^2/2-i\gamma t^3/6)$. Taylor's integral remainder
for $e^{itX}$, followed by expectation, gives its expansion through
degree five with error at most $\mu_6t^6/720$. For
$w=-t^2/2-i\gamma t^3/6$ and $0\le t\le1/2$,
\[
 |w|\le t^2(1/2+g^2/144)<1/6,
 \quad |e^w-1-w-w^2/2|\le |w|^3e^{|w|}/6\le |w|^3/5.
\]
The first bound follows from $\sqrt{1/4+y}\le1/2+y$; for the last
use $e^{1/6}\le(1-1/6)^{-1}=6/5$. Comparing the two polynomials gives
\begin{equation}\label{j:fourier:taylor}
 f-h=-\mathfrak a t^4-i\mathfrak z t^5+r_6(t),\qquad |r_6(t)|\le Et^6,
 \qquad |f-h|\le Bt^4.
\end{equation}
The degree-six term in $w^2/2$ costs $g^2t^6/72$; it has not been
discarded.

Four integrations of $1-\cos u\ge0$ show
$\cos u\le1-u^2/2+u^4/24$ for every real $u$. Independence gives
$\E(X-X')^4=2\mu_4+6$, and hence
\[
 |f(t)|^2\le1-t^2+(M+3)t^4/12
 \le1-(45-M)t^2/48.
\]
Thus $|f|,|h|\le e^{-(45-M)t^2/96}$ throughout the low interval.
Every point on their complex joining segment obeys the same bound.
Taylor's integral formula for the polynomial $v\mapsto v^n$ proves
\[
 f^n-h^n=n(f-h)h^{n-1}+R_2,\quad
 |R_2|\le\tfrac12n(n-1)B^2t^8e^{-(n-2)(45-M)t^2/96}.
\]
The leading remainder in \eqref{j:fourier:taylor} and this power
remainder cost at most $L_2/n$ after division by $t^3$ and integration:
use $n-1\ge2a_0n$, $(n-2)(45-M)/96\ge b_0n$, $n(n-1)\le n^2$,
and $\int_0^\infty t^3e^{-At^2}dt=1/(2A^2)$,
$\int_0^\infty t^5e^{-At^2}dt=A^{-3}$.

For the remaining leading terms let $b=n-1$ and
$v=b\gamma t^3/6$. Their imaginary part after rotation by $Ht$ is
bounded by
\[
 ne^{-bt^2/2}\{(Z+A|H|)t^5+bAg t^7/6\}.
\]
The Gaussian moments of orders two and four, after dividing by $t^3$,
give $n\sqrt{\pi/2}b^{-3/2}(Z+A|H|+Ag/2)$.
This is at most $\lambda_N(Z+A|H|+Ag/2)/\sqrt n$.
Finally put $u=n\gamma t^3/6$. Direct sine and cosine remainders give
\[
 |\Im[e^{-itH}(h^n-e_n)]|
 \le e^{-nt^2/2}(|u|^3/6+|H|t u^2/2).
\]
The Gaussian moments of orders six and four respectively give
$\sqrt{\pi/2}g^3/(432\sqrt n)$ times $5$, and
$\sqrt{\pi/2}g^2|H|/(24\sqrt n)$.
Using $\sqrt{\pi/2}<251/200$ gives precisely the two displayed
rational coefficients. Every extended integrand was nonnegative.
The only index comparisons were the displayed rational inequalities;
in particular the law and $H$ can vary arbitrarily with $n$.
\end{proof}

\paragraph{A common Pearson envelope.}\label{j:fourier:pearson-envelope}
For actual support $|X|\le R\le6$ and $\E|X|^5\le E_5$, the
triangle LOW bound may use $M_6=RE_5$ and $A=\max(2,M-3)/24$.
To bound its fifth coefficient, put $Q=X^2-\gamma X-1$. Multiplication
and the first three standardized moments give
\[
 \E Q^2=\mu_4-1-\gamma^2,\qquad
 \E X^5=\gamma^3+2\gamma+\E[Q^2(X+2\gamma)].
\]
Reflect this auxiliary calculation so that $t=|\gamma|\ge0$.
For $r\ge\max(R,2g)$ the two signs of $\E X^5-10\gamma$ are at most
\[
 P_r(t)=r(M-1)+(10-2M)t-rt^2+t^3,
 \qquad Q_r(t)=r(M-1)+(2M-10)t-rt^2-t^3.
\]
Indeed the coefficients of $\mu_4$ before replacing it by $M$ are
$r-2t,r+2t\ge0$. More precisely, for
$z=(10\gamma-\E X^5)/120$ this proves $-Q_r(t)\le120z\le P_r(t)$.
If $M\le5$, $P_r\ge Q_r$ and
$P_r\le r(M-1)+(10-2M)t-(r-g)t^2$.
If $M>5$, $P_r\le r(M-1)$ and
$Q_r\le r(M-1)+(2M-10)t-rt^2$.
Both quadratics are maximized on $[0,g]$ at their clipped vertices.
Thus for $r_0=\max(R,2g)>g$ the exact selected value is
\begin{gather*}
 Z_r=r_0(M-1)+a t_*-b t_*^2,\qquad
 t_*=\min\{g,a/(2b)\},\\
 (a,b)=\begin{cases}(10-2M,r_0-g),&M\le5,\\
                    (2M-10,r_0),&M>5.
       \end{cases}
\end{gather*}
For the radius-six choice $P_6''=-12+6t<0$ on $[0,g]$.
Consequently its selected ceiling is
\[
 Z_0=\begin{cases}
 P_6(g),&M\le5,\ P'_6(g)\ge0,\\
 6(M-1)+(10-2M)^2/[4(6-g)],&M\le5,\ P'_6(g)<0,\\
 6(M-1)+(2M-10)^2/24,&M>5.
 \end{cases}
\]
The first branch uses monotonicity up to $g$; the others use the
unrestricted quadratic vertices just proved. Hence
$120Z=\min\{E_5+10g,Z_0,Z_r\}$ bounds the actual fifth coefficient.
Every radius, moment and sign assumption refers to this same marginal.

\subsection{The complete high-frequency half-line}

A nearly lattice law can retain characteristic-function peaks away from
zero. The high-frequency argument therefore uses an event involving the
actual difference $X-X'$ and partitions the oscillations around their
zeros. The six contributions below cover intermediate frequencies,
nonzero-cell cores, nonzero-cell wings, the small-angle part of the zero
cell, the rest of the zero cell, and the signed Gaussian comparator.
Each optional improvement has its own stated positivity and monotonicity
conditions.

For independent copies of the same marginal suppose an actual event
$\mathcal E$ satisfies
\[
 \Pp(\mathcal E)\ge\alpha>0,\quad |X-X'|\ge d>0\text{ on }\mathcal E,
 \quad\E[(X-X')^2\mid\mathcal E]\le M_E.
\]
Fix $b=9/10$, $0<r\le1/4$, and
\[
 h=333/106-(22/7)^3r^2/6>0,\quad
 \kappa=\alpha(23/48)^2d^2,\quad L(x)=\frac{x+2}{x^2+4x+2},
\]
\[
 S_3=\sum_{j=1}^{20}j^{-3}+1/800,\quad
 S_5=\sum_{j=1}^{20}(j-r)^{-5}+\frac1{4(20-r)^4}.
\]
Let $c(u)=u/2-(M+3)u^2/24$ and set
\[
 c_* =\min\{c(1/4),c(81/100)\},\quad
 c_q=1/8-(M+3)/384,\quad s=[1-(M+3)53/600]/2.
\]
\begin{proposition}\label{j:fourier:high}
Suppose $c_*>0$. For every $n\ge N$, $\sqrt n I_\infty$ is at most
the sum of the following six terms:
\begin{align*}
 T_1(n)&=2\sqrt n e^{-nc_*},\\
 T_2(n)&=\frac{M_E}{4\pi^{3/2}h\sqrt\alpha}
       \{S_3+3S_5/(n\alpha h^2)\},\\
 T_3(n)&=\frac{M_E\sqrt n}{2\pi^2}e^{-n\alpha h^2r^2},\\
 T_4(n)&=\frac{e^{-n\kappa b^2}}{2b^4\kappa\sqrt n},\\
 T_5(n)&=\frac{\sqrt n}{2b^2}e^{-n\alpha(23/48)^2},\\
 T_6(n)&=\sqrt n e^{-n/8}\{2L(n/8)+\sqrt{M-1}/3\}.
\end{align*}
If $s>0$, $T_1$ may be replaced by its minimum with
$2\sqrt n e^{-nc_q}L(ns/4)$. The fourth term may be replaced by
its minimum with $\sqrt n e^{-n\kappa b^2}L(n\kappa b^2)/(2b^2)$.
The common-floor sum at $N$ is valid if
\[
 2Nc'>1\quad\text{for }c'\in\{c_*,\alpha h^2r^2,
                      \alpha(23/48)^2,1/8\}.
\]
For the optional first term also require $2Nc_q>1$; for the optional
fourth term require $N\kappa b^2\ge1/4$. If an optional gate fails,
the simpler term remains available.
\end{proposition}
\begin{proof}
The triangle inequality separates $|f^n|$ and the entire comparator.
On $[1/2,b]$, symmetrization gives $|f|^n\le e^{-nc(t^2)}$.
The function $c$ is concave, so $c(t^2)\ge c_*$ and the integral
of $t^{-3}$ is less than two. This gives $T_1$.
Also
\[
 c(u)-c_q-s(u-1/4)=(M+3)(u-1/4)(81/100-u)/24\ge0.
\]
When $s>0$ this chord may be integrated to infinity.

For completeness, the tail calculation used here is, for $x=na b^2>0$,
\begin{equation}\label{j:fourier:tail}
 \sqrt n\int_b^\infty e^{-na t^2}t^{-3}dt
 \le\frac{\sqrt n e^{-x}}{2b^2}L(x).
\end{equation}
To see where the rational function $L$ comes from, let $U$ have density
$e^{-u}$ on $u\ge0$. Substitution followed by positive Fubini gives
\[
 \sqrt n\int_b^\infty e^{-na t^2}t^{-3}dt
 =\frac{\sqrt n e^{-x}}{2b^2}\E\frac{U}{x+U}.
\]
For example, after putting $v=na(t^2-b^2)$ the remaining factor is
$\int_0^\infty xe^{-v}(x+v)^{-2}dv$; integration by parts makes it
$\E[U/(x+U)]$. For $x>0$ put
\[
 J_x=\E\frac1{x+U},\qquad v_0=\begin{pmatrix}1\\1\end{pmatrix},\qquad
 G=\E\left[(x+U)\begin{pmatrix}1&U\\U&U^2\end{pmatrix}\right]
 =\begin{pmatrix}x+1&x+2\\x+2&2x+6\end{pmatrix}.
\]
The entries use $\E U=1$, $\E U^2=2$, and $\E U^3=6$.
For $q(U)=c_1+c_2U$, weighted Cauchy--Schwarz in its squared-residual
form gives
\[
 0\le\E\frac{[1-(x+U)q(U)]^2}{x+U}
 =J_x-2v_0^Tc+c^TGc.
\]
The matrix $G$ is positive definite: its leading entry is positive and
$\det G=x^2+4x+2>0$. Choosing $c=G^{-1}v_0$ yields
\[
 J_x\ge v_0^TG^{-1}v_0=\frac{x+3}{x^2+4x+2},\qquad
 \E\frac{U}{x+U}=1-xJ_x\le\frac{x+2}{x^2+4x+2}=L(x).
\]
This proves \eqref{j:fourier:tail} for every $a,b,n>0$ with the same
function $L$. Applying it at endpoint
$1/2$ and rate $s$ proves the optional first term.

For $t\ge b$, symmetrization, the actual-event probability bound,
and convex Jensen give
\[
 |f(t)|^n\le e^{-n\E\sin^2(t(X-X')/2)}
 \le\E[e^{-n\alpha\sin^2(tD/2)}\mid\mathcal E],
 \quad D=|X-X'|.
\]
Positive Tonelli allows a separate partition for every $D>0$.
With $\ell=2\pi/D$, use the zero cell $[0,\ell/2]$ and nonzero
cells $\ell[j-1/2,j+1/2]$, $j\ge1$. In a nonzero cell write
$t=\ell(j+v)$. On its core $|v|\le r$,
$|\sin\pi v|\ge h|v|$; on its wings $r\le|v|\le1/2$ it is
at least $hr$. These follow from $\sin x\ge x-x^3/6$ and
monotonicity of $\sin$ on $[0,\pi/2]$, with directed bounds on $\pi$.
For $|v|\le r$ Taylor's integral remainder gives
\[
 \{(j-v)^{-3}+(j+v)^{-3}\}/2
 \le j^{-3}+6(j-r)^{-5}v^2.
\]
Extend the nonnegative Gaussian majorant over the whole real line.
Its zeroth and second moments, the Jacobian
$t^{-3}dt=D^2(4\pi^2)^{-1}(j+v)^{-3}dv$, and
$\E[D^2\mid\mathcal E]\le M_E$ give $T_2$.
The two series tails are bounded by their decreasing integrals after
20. The total wing weight is at most
$\int_{\ell/2}^\infty t^{-3}dt=D^2/(2\pi^2)$, giving $T_3$.
Extending these sets below $b$ only increases their nonnegative bounds.

In the zero cell put $u=tD/2\le\pi/2$. Where $u\le1/2$,
$\sin u\ge23u/24$, and $D\ge d$ yields rate $\kappa$.
Extending the set to $[b,\infty)$ gives $T_4$ from
$t^{-3}\le t/b^4$; \eqref{j:fourier:tail} gives its optional improvement.
Where $u\ge1/2$, monotonicity gives $\sin u\ge23/48$ and the total
weight is at most $1/(2b^2)$, proving $T_5$. Endpoints have zero
Lebesgue measure even if the conditional distance distribution has atoms.

Finally $|e_n(t)|\le e^{-nt^2/2}(1+n|\gamma|t^3/6)$.
The first part on the entire $[1/2,\infty)$ uses
\eqref{j:fourier:tail}; the second uses
$\int_{1/2}^\infty e^{-nt^2/2}dt\le2e^{-n/8}/n$.
Pearson's identity $\mu_4-1-\gamma^2=\E(X^2-\gamma X-1)^2\ge0$
gives $T_6$. No comparator tail is dropped.

For floor evaluation, $\sqrt n e^{-nc'}$ decreases for $2nc'>1$,
whereas $n^{-1}$ and $n^{-1/2}e^{-nc'}$ decrease for every $c'>0$.
Also $L'(x)=-(x^2+4x+6)/(x^2+4x+2)^2<0$.
Thus the optional first term decreases under its own gate, and the
comparator decreases for $n\ge5$. The sign of the derivative of
$\sqrt{x}e^{-x}L(x)$ is the sign of
\[
 4-10x-22x^2-13x^3-2x^4,
\]
which is negative at $1/4$ and strictly decreasing on $(0,\infty)$.
This proves the optional zero-cell gate. Every comparison is of a common
scalar majorant, not of characteristic functions for different laws.
\end{proof}

\subsection{Uniform remainder conclusion}

Under the preceding moment, actual-event and scalar gates, let $K_H$ be
any complete HIGH sum evaluated at $N$, including its permitted optional
improvements. Then
\begin{equation}\label{j:fourier:conclusion}
 |E_n|\le\frac2\pi\{L_0+K_H+L_2/\sqrt N+L_H|H|\}
 \le\frac{16}{25}\{L_0+K_H+L_2/\sqrt N+L_H|H|\}.
\end{equation}
The second inequality uses $\pi>25/8$ and nonnegative coefficients.
This supplies a constant Fourier budget and a raw-threshold slope.
The scalar gates are imposed on the analytic majorants, not on the rounded
implementation. Complete HIGH bounds for different valid windows may be
minimized. Each LOW triple remains intact: a constant from one enclosure
cannot be combined with another enclosure's threshold slope without a
separate justification. These statements are uniform over every law and
threshold satisfying the displayed hypotheses; no fixed-law limit is used.
